% Figure: bar chart comparing the stored energy per cubic metre for several
% thermal-storage media and modes (hot water, concrete, molten salt sensible,
% PCM latent), giving a sense of energy density across the three modes.
\begin{figure}[htbp]
\centering
\begin{tikzpicture}
\begin{axis}[
  width=12cm, height=7.5cm,
  ybar,
  bar width=18pt,
  xlabel={storage medium and mode},
  ylabel={stored energy density (kWh\,m$^{-3}$)},
  ymin=0, ymax=260,
  symbolic x coords={water,concrete,salt,PCM,thermochem},
  xtick=data,
  x tick label style={font=\scriptsize, rotate=20, anchor=north east},
  tick label style={font=\scriptsize},
  label style={font=\small},
  nodes near coords,
  every node near coord/.append style={font=\scriptsize},
  enlarge x limits=0.15,
]
% values: hot water dT=60 ~ 70; concrete ~ 24; molten salt sensible ~ 130;
% PCM latent ~ 50; thermochemical (research) ~ 200 (high but immature)
\addplot[fill=engblue!70] coordinates {
  (water,70) (concrete,24) (salt,130) (PCM,50) (thermochem,200)
};
\end{axis}
\end{tikzpicture}
\caption[Thermal-storage energy density by medium and mode]{Approximate stored
energy density for several thermal-storage media and modes. Hot water and
concrete store sensible heat over a usable temperature swing; molten salt stores
sensible heat too, but over a much larger swing and so reaches a higher density;
the phase-change material stores latent heat at a fixed melting temperature; the
thermochemical bar is the highest but rests on systems still mostly in research as
of 2026, so the bar reflects potential rather than installed practice. The figure
fixes the order of magnitude that sets the tank size for a given storage duty:
sensible water storage needs roughly three times the volume of a molten-salt
store for the same energy, which is why the medium is chosen against the
available temperature range, not the density alone.}
\label{fig:vol04ch06:storage-comparison}
\end{figure}
